\section{Conclusiones y recomendaciones}

En primer lugar, respecto a las conclusiones del proyecto realizado, estas se enumeran a continuación:

\begin{itemize}
    \item Se logró una comprensión sólida de la cláusula 36 del estándar IEEE~802.3, lo cual permitió definir el comportamiento esperado de la PCS y de cada uno de los módulos desarrollados.
    \item Se consiguió una implementación adecuada con el funcionamiento esperado de uno de los módulos de interés que integran la PCS, así como un resultado satisfactorio para la integración de los tres.
    \item El plan de pruebas desarrollado permitió asegurar el funcionamiento, tanto en el caso idóneo como en algunos casos de error (especialmente en el caso del sincronizador). De esta forma, la integración de los tres resultó más sencilla, pues ya se conocía el comportamiento individual de los módulos.
    \item El uso de un enfoque modular resultó adecuado y eficiente, permitiendo que cada integrante del grupo trabajara en un bloque independiente sin comprometer la coherencia del diseño total.
    \item La adopción de convenciones unificadas en nombres de señales, estilos de codificación y estructura de archivos facilitó considerablemente la revisión del código y la colaboración entre los participantes.
    \item El repositorio y flujo de trabajo en GitHub favorecieron una gestión ordenada del proyecto, evitando conflictos y permitiendo avanzar de manera paralela en diferentes componentes.
\end{itemize}

En cuanto a las recomendaciones para la siguiente etapa del proyecto, se proponen las siguientes:

\begin{itemize}
    \item Realizar la síntesis de los módulos.
    \item Implementar casos de configuración y de error en cada uno de los submódulos para tener un desarrollo completo de la cláusula del estándar.
    \item Explorar la implementación del bloque de Auto-Negotiation y Carrier-Sense para completar la PCS del estándar IEEE 802.3.
\end{itemize}
